\documentclass[11pt]{article}
\usepackage[bottom]{footmisc} 
\renewcommand{\thempfootnote}{\arabic{mpfootnote}}

\usepackage[hidelinks]{hyperref}

\usepackage[utf8]{inputenc}

\usepackage[T1]{fontenc}
\usepackage[ttdefault=true]{AnonymousPro}

\usepackage{multicol}
\setlength{\columnsep}{1cm}

\usepackage{abc}
\usepackage{graphicx}
\usepackage{musicography}

\usepackage[margin=0.5in]{geometry} 
\geometry{letterpaper}


\usepackage{fancyhdr} 
\pagestyle{fancy} % options: empty , plain , fancy
\renewcommand{\headrulewidth}{0pt} % customise the layout...
\lhead{}\chead{}\rhead{}
\lfoot{}\cfoot{\thepage}\rfoot{}


\begin{document}

\thispagestyle{empty}

\begin{abc}[name=girlileftbehind,options={-O= +c -F single.fmt}]
X:1
T:The Girl I Left Behind Me
M:2/4
L:1/8
K:G
|:g/f/ | "G"ed BG | "C"AG ED | "G"GG/F/ G/A/B/c/ | "D"d2 Bg/f/ | "G"ed BG | "C"AG EG | "D"FA DE/F/ | "G"G2 G :|
|: G/A/ | "G"Bd ef | gd BG | Bd ef | "D"g2 fg/f/ |  "G"ed BG | "C"AG EG | "D"FA DE/F/ | "G"G2-G :|
\end{abc}

\vskip 0.5in
\noindent\makebox[\textwidth][c]{%
  \begin{minipage}{7.0in}
    \setlength{\parindent}{1.5em}

    Both England and Ireland claim to have originated The Girl I Left
    Behind Me. The tune appears as The Spailpin Fanach (The Rambling
    Laborer) in Dublin in 1791 although some claim it was around much
    earlier. In England, the tune was a marching air from around the
    beginning of the 19th century. Edward Bunting included in a letter
    in 1840, ``It is a pretty tune, and has been played for the last
    fifty years, to my knowledge, by the fifes and drums, and bands of
    different regiments, on their leaving the towns for new quarters.''\footnote{\url{https://tunearch.org/wiki/Annotation:Girl\_I\_Left\_Behind\_Me\_(1)\_(The)}}

    
    John Bell, a Northumbrian (North East England) musician, poet, and
    writer, includes the below lyric for Blyth Camps (or The Girl I
    Left Behind Me) in his collection, \textbf{Rhymes of the Northern Bards}
    (Newcastle upon Tyne, 1812):\footnote{\textit{Loc. cit.}}

    {\parindent0pt
      \raggedright
      ~\\
      \textit{
      I'm lonesome since I left Blyth Camps, \\
      And o'er the moor that's sedgy; \\
      With heavy thoughts my mind is fill'd \\
      Since I parted with my Betsy, \\
      Whene'er I turn to view the place \\
      The tears fall down and blind me \\
      When I think on the charming grace \\
      Of her I left behind me.--- \\
      }
      ~\\
    }
    
    In America, legend has it that it was overheard being sung by a
    British prisoner of war during the War of 1812 and caught on among the
    American infantries. It was played by both the Union and Confederacy
    armies during the American Civil War and was a favorite tune of
    General George Custer.\footnote{\textit{Loc. cit.}}
    
    Eloise Hubbard Linscott includes a The Girl I Left Behind Me
    contradance in her collection, \textbf{Folk Songs of Old New England} (Hamden
    [CT], 1939):\footnote{Linscott. \textbf{Folk Songs of Old New England}, 1939, p 79.}

    {\parindent0pt
      \raggedright
      ~\\
      \begin{tabular}{ll}
        Circle eight hands around & 8 bars \\
        First couple lead to the right halfway round & 4 bars \\
        Pass right through between the two & 4 bars \\
        Swing the girl behind you & 8 bars \\
        Ladies grand chain & 8 bars \\
        All promenade with the dear little maid, that girl you left behind you & 8 bars \\
      \end{tabular}
      ~\\
    }
    ~\\
    ---Kitty Steetle
    
  \end{minipage}
  \vspace*{\fill}
}

\end{document}
